\documentclass[11pt, reqno]{amsart}

\usepackage[spanish]{babel}
\usepackage[LGR, T1]{fontenc}
\usepackage[utf8]{inputenc}

../LaTeX/general.tex
% \input{../graphics.tex}

\makeatletter
\def\emailaddrname{\textit{Correo electrónico}}
\def\subtitle#1{\gdef\@subtitle{#1}}
\def\@subtitle{}

% Metadata
\def\logo#1{\gdef\@logo{#1}}
\def\@logo{}
\def\institution#1{\gdef\@institution{#1}}
\def\@institution{}
\def\department#1{\gdef\@department{#1}}
\def\@department{}
\def\professor#1{\gdef\@professor{#1}}
\def\@professor{}
\def\course#1{\gdef\@course{#1}}
\def\@course{}
\def\coursecode#1{\gdef\@coursecode{#1}}
\def\@coursecode{}

\renewcommand{\maketitle}{
\begin{center}
	\small
	\renewcommand{\arraystretch}{1.2}
	\begin{tabular}{cp{.37\textwidth}p{0.44\textwidth}}
		% \hline
		\multirow{5}{*}{\includegraphics[height=2.0cm]{\@logo}}
	  & \multicolumn{2}{c}{ \makecell{{\bfseries \@institution} \\ \@department} } \\
	  % & \multicolumn{2}{c|}{{\bfseries\@institution} \\ \@department} \\
	  \cline{2-3}
	  & \textbf{Profesor:} \@professor & \textbf{Ayudante:} \authors \\
	  % \cline{2-3}
	  & \textbf{Curso:} \@course & \textbf{Sigla:} \@coursecode \\
	  % \cline{2-3}
	  & \multicolumn{2}{l}{ \textbf{Fecha:} \@date } \\
	  % \hline
	\end{tabular}
	\\[\baselineskip]
	% {}
	% \vspace{2\baselineskip}
	{\bfseries\Large\@title}
	\ifx\@subtitle\@empty\else
		\\[1ex]
		\large\mdseries\@subtitle
	\fi
\end{center}
}
\makeatother

\usepackage{multirow, makecell}

\usepackage[
	reversemp,
	letterpaper,
	% marginpar=2cm,
	% marginsep=1pt,
	margin=2.3cm
]{geometry}
\usepackage{fontawesome}
% \makeatletter
% \@reversemargintrue
% \makeatother

% Símbolos al margen, necesitan doble compilación
\newcommand{\hard}{\marginnote{\faFire}}
\newcommand{\hhard}{\marginnote{\faFire\faFire}}

% Dependencias para los teoremas
\usepackage{xifthen}
\def\@thmdep{}
\newcommand{\thmdep}[1]{
	\ifthenelse{\isempty{#1}}
	{\def\@thmdep{}}
	{\def\@thmdep{ (#1)}}
}
\newcommand{\thmstyle}{\color{thm}\sffamily\bfseries}

% ===== Estilos de Teoremas ==========
\newtheoremstyle{axiomstyle}
	{0.3cm}
	{0.3cm}
	{\normalfont}
	{0.5cm}
	{\bfseries\scshape}
	{:}
	{4pt}
	{\thmname{#1}\thmnote{ #3}\thmnumber{ (#2)}}
\newtheoremstyle{styleC}
	{0.5cm}
	{0.5cm}
	{\normalfont}
	{0.5cm}
	{\bfseries}
	{:}
	{4pt}
	{\thmname{#1\textrm{\@thmdep}}\thmnumber{ #2}\thmnote{ (#3)}}

% ====== Teoremas (sin borde) ===========
\theoremstyle{axiomstyle}
\newtheorem*{axiom}{Axioma}

% ====== Teoremas (sin borde) ==================
\theoremstyle{styleC}
\newtheorem{thm}{Teorema}[section]
\newtheorem{mydef}[thm]{Definición}
\newtheorem{prop}[thm]{Proposición}
\newtheorem{cor}[thm]{Corolario}
\newtheorem{lem}[thm]{Lema}
\newtheorem{con}[thm]{Conjetura}
\newtheorem*{sol}{Solución}
\newtheorem*{obs}{Observación}
\newtheorem*{ex}{Ejemplo}

\newtcbox{bluebox}[1][]{enhanced jigsaw, 
  sharp corners,
  frame hidden,
  nobeforeafter,
  listing only,
  #1} % comando para crear cajas de colores

\expandafter\let\expandafter\oldproof\csname\string\proof\endcsname
\let\oldendproof\endproof
\renewenvironment{proof}[1][\proofname]{%
  \oldproof[\scshape Demostración:]%
}{\oldendproof} % comando para redefinir la caja de la demostración
\newenvironment{hint}[1][\proofname]{%
  \oldproof[\scshape Pista:]%
}{\oldendproof} % comando para redefinir la caja de la demostración

% colores utilizados
\definecolor{numchap}{RGB}{249,133,29}
\definecolor{chap}{RGB}{6,129,204}
\definecolor{sec}{RGB}{204,0,0}
\definecolor{thm}{RGB}{106,176,240}
\definecolor{thmB}{RGB}{32,31,31}
\definecolor{part}{RGB}{212,66,66}

% ====== Diseño de los titulares ===============
\usepackage[explicit]{titlesec} % para personalizar el documento, la opción <<explicit>> hace que el texto de los titulares sea un objeto interactuable

\titleformat{\subsection}[runin]
	{\bfseries}
	{\textrm{\S}\thesubsection}
	{1ex}
	{#1.}

\setlist[enumerate,1]{label=\arabic*., ref=\arabic*} % Enumerate standards

\usepackage{tikz}
\usetikzlibrary{babel,cd}
% \DeclareMathOperator\Gr{Gr}

\title{El teorema de coeficientes universales}
\date{\DTMdate{2026-08-20}}

\author{José Cuevas Barrientos}
\email{josecuevasbtos@uc.cl}
\urladdr{https://josecuevas.xyz/teach/2026-2-ayud/}

\logo{../puc_negro.png}
\institution{Pontificia Universidad Católica de Chile}
\department{Facultad de Matemáticas}
\course{Topología Algebraica}
\coursecode{MAT2850}
\professor{Mauricio Bustamante}

\begin{document}

\maketitle

\section{Cohomología}
\begin{mydef}
	Un \emph{cocomplejo} de módulos, denotado $A^\bullet$, será una familia $(A^n)_{n\in\Z}$ de módulos con diferenciales $\ud^n_A \colon
	A^n \to A^{n+1}$ \textquote{creciendo en grado} tales que $\ud_A^{n+1}\circ\ud_A^n = 0$ para todo $n \in \Z$.
\end{mydef}
Uno puede imitar la definición de homomorfismo de complejos para cocomplejos y de homología para obtener la \emph{cohomología}:
\[
	H^n(A) := \frac{\ker( \ud \colon A^n \to A^{n+1} )}{\Img( \ud \colon A^{n-1} \to A^n )}.
\]
Por ejemplo, es inmediato que:
\begin{obs}
	Las categorías de complejos y de cocomplejos son isomorfas mediante la aplicación que a un complejo $A_\bullet$ le asocia el
	cocomplejo $A^n := A_{-n}$ con $\ud^n_A := \ud^A_n$.
	Más aún, $H_n(A_\bullet) = H^n(A^\bullet)$ (aquí hay honesta igualdad).
\end{obs}

\begin{mydef}
	Sea $A$ un $R$-módulo.
	Una \strong{resolución libre} es un complejo de módulos exacto
	\[\begin{tikzcd}
		\cdots \rar & F_2 \rar & F_1 \rar & F_0 \rar["\varepsilon"] & A
	\end{tikzcd}\]
	donde cada $F_n$ es un módulo libre.
	Denotaremos por $F_\bullet$ al complejo (sin incluír a la flecha $\varepsilon \colon F_0 \epicto A$).
\end{mydef}

\begin{enumerate}
	\item\lookright
		Pruebe que todo $R$-módulo admite una resolución libre.

	\item Sea $A_\bullet$ un complejo de módulos.
		\begin{enumerate}
			\item Muestre que para todo $R$-módulo $B$, la familia $C^n := \Hom_R(A_n, B)$ con las aplicaciones
				\[
					\ud^n_C \colon C^n \longrightarrow C^{n+1}, \qquad \varphi \longmapsto \varphi\circ\ud^A_n
				\]
				forma un cocomplejo.

			\item\lookup ¿Se cumple que $H^n(C^\bullet) \cong H_n(A_\bullet)$?
		\end{enumerate}
\end{enumerate}

\section{El funtor $\Ext$}
Puede tomar por sabido el siguiente hecho:
\begin{prop}[{\cite[129-131]{hilton:homological}}]
	Sean $A, B$ un par de $R$-módulos.
	Se cumple que
	\[
		\forall n\in\Z, \qquad
		H^n( \Hom(F_\bullet, B) ) \cong H^n( \Hom(G_\bullet, B) ),
	\]
	donde $F_\bullet \to A$ y $G_\bullet \to A$ son dos resoluciones libres.
	A este módulo lo denotamos por $\Ext_R^n(A, B)$.
\end{prop}
\begin{enumerate}[resume]
	\item Pruebe que $\Ext_R^0(A, B) \cong \Hom_R(A, B)$.

	\item\lookright Calcule $\Ext^1_\Z(\Z, \Z/n\Z)$ y $\Ext^1_\Z(\Z/n\Z, \Z)$.

	\item Sea $0 \to A \to B \to C \to 0$ una sucesión exacta de módulos, pruebe que para todo módulo $M$ hay una sucesión exacta
		\[\begin{tikzcd}
			0 \rar & \Hom(C, M) \rar & \Hom(B, M) \rar & \Hom(A, M) \\
			{} & \Ext^1(C, M) \rar & \Ext^1(B, M) \rar & \Ext^1(A, M)
			\arrow["\partial"{description}, from=1-4, to=2-2, out=-150, in=40]
		\end{tikzcd}\]
		% \begin{hint}
		% 	Neces
		% \end{hint}

	\item\lookst \textbf{Teorema de coeficientes universales:}
		Pruebe que para todo complejo $A_\bullet$ de \underline{grupos abelianos} \underline{libres} y todo grupo abeliano $M$, tenemos una sucesión exacta
		\[\begin{tikzcd}
			0 \rar & \Ext^1(H_{n-1}(A_\bullet), M) \rar & H^n(\Hom(A_\bullet, M)) \rar & \Hom(H_n(A_\bullet), M) \rar & 0.
		\end{tikzcd}\]
		\begin{hint}
			% Dada una resolución libre $F_\bullet \to A$, note que los núcleos $K_n := \ker(\ud \colon F_n \to F_{n-1})$ son
			% también libres.
			Podrían resultar útiles los siguientes datos:
			\begin{itemize}
				\item Todo subgrupo de uno abeliano libre es libre.
				\item Si $K_n := \ker(\ud \colon B_n \to B_{n-1})$ denota el núcleo de un complejo arbitrario $B_\bullet$, tenemos una sucesión exacta de
					complejos $0 \to K_\bullet \to B_\bullet \to \ud B_\bullet \to 0$.
				\item Y para cada $n\in\Z$ una sucesión exacta $0 \to \ud B_{n+1} \to K_n \to H_n(B) \to 0$ de grupos.
					\qedhere
			\end{itemize}
		\end{hint}
\end{enumerate}

\section*{Comentarios adicionales}
El nombre \textquote{Ext} tiene que ver con que $\Ext^1(A, B)$ clasifica \textquote{extensiones} $C$ de $B$ sobre $A$, es decir, objetos que
calzan en una sucesión exacta corta $0 \to B \to C \to A \to 0$.
De hecho, uno podría tomar esto como definición para obtener un conjunto, ocasionalmente llamado el \textquote{funtor Ext de Yoneda}, al
cual podemos dotar manualmente de una operación de grupo (cfr.\ \cite{hilton:homological}, \S IV.9); este enfoque tiene sus ventajas al
trabajar con grupos no conmutativos.

Se entienden por \textquote{teorema de coeficientes universales} a otros enunciados similares, entre ellos, uno que dice cómo calcular la
homología del tensor de un complejo $A_\bullet \otimes_R S$ con \textquote{otros coeficientes} (cfr.\ \cite{hilton:homological}, Th.~V.2.5);
de ahí el nombre.

El \textquote{teorema de coeficientes universales en \emph{cohomología}} también es cierto cuando $A_\bullet$ y $M$ son módulos sobre un DIP
sin ninguna modificación.
Si el anillo base no es un DIP, existen submódulos de uno libre que no son libres (¿se le ocurre un ejemplo?) y el estudiante puede
verificar que el argumento se cae.

\printbibliography

\end{document}
